\scanpage{171}
% L1–4
\[
\begin{aligned}
x_1*x_2&=(o_1+x_1)*(\unsure{o|o_2}+x_2)
=o_1*o_2+\bigl(x_1\otimes(\overbrace{\sigma_2(o_2)}^{b_2}-o_2\bigr)\\
&\qquad+(\sigma_1(o_1)-o_1)\otimes x_2-2\deleted{\unsure{f}}x_1\otimes x_2\\
&=o+\underbrace{x_1\otimes b_2+b_1\otimes x_2-2x_1x_2}\\
&=x_1b_2+b_1x_2-2x_1x_2
\end{aligned}
\]
% L5–10
$\pi_1:A_1\longrightarrow Q$ transforme $x_1$ en la forme
\[
\pi_1(x_1)=q_{x_1(b-x_1)}
=\bigl(z_1\mapsto {\deleted{\unsure{x}}z}_1^2-b_1z_1+
\deleted{\unsure{(}}x_1(b_1-x_1)\unsure{'}=(z_1-x_1)(z_1-(b_1-x_1))\bigr)
\]
et de $=$ pour $\pi_2$, donc
$\pi_1(x_1)\otimes\pi_2(x_2)$ est la forme
\[
\begin{aligned}
z\longmapsto{}&z^2-b_1b_2z+b_2^2x_1(b_1-x_1)+b_1^2x_2(b_2-x_2)\\
&-4x_1(b_1-x_1)x_2(b_2-x_2)
\end{aligned}
\]
% L11–12
et
\[
\pi(x_1*x_2)=\bigl(z\mapsto z^2-\underbrace{b_1b_2}_{b}z+
\underbrace{(x_1*x_2)(b_1b_2-x_1*x_2)}_{\substack{(x_1b_2+b_1x_2-2x_1x_2)(b_1b_2-x_1b_2\unsure{+|-}b_1x_2\\{}+2x_1x_2)}}
\]

% L13–19
Le calcul de l'égalité est un peu pénible
mais ça peut se faire (j'ai dû le
vérifier précéd$^{\mathrm{mt}}$ par un autre biais,
en choisissant $q_1\,q_2$ et faisant la vérification
en choisissant $o_i=x_i$ \deleted{tel que $q_{\unsure{i|1}}(o_i)=0$} i.e
il suffit de le faire pour $\unsure{x_1'|x_1}=0$, $\unsure{\alpha_2|x_2}=0$
où \unsure{c'est|est} trivial.

% L20–26
On voudrait écrire \deleted{$Q(A)=Q(A$}
\[
Q(A_1*A_2)\simeq Q(A_1)*Q(A_2)
\]
mais la théorie du $*$ développée
précédemment ne s'y prête pas, \unsure{car},
\deleted{$Q$} les $Q(A_i)$ ne sont pas involutifs.
Il faut reprendre \deleted{tout} le \unsure{formalisme|formulation|formalisation} \inserted{\unsure{$*$|$A$|($A$}}
de façon plus générale!
% archive sheet number: 170
